\documentclass[12pt]{article}
\usepackage{amsmath}
\usepackage{geometry}% 1-inch margins
\geometry{letterpaper, margin=1in}% 1-inch margins

\title{Mass-to-Force Conversions in English Units}
\author{Jacob C. Vaught}
\date{Fall 2023 EMCH 290}

\begin{document}
\maketitle

\section*{Introduction}
This document aims to provide a brief look into three english[USCS] units: pound-mass (lbm), pound-force (lbf), and slug. The definitions will be broken down mathematically and examples will be provided to illustrate their conversions and applications.

\subsection*{Definitions}

\subsubsection*{Pound-Force (lbf)}
The pound-force (\( \text{lbf} \)) is an English unit used to quantify force. By definition, it is the force required to accelerate a one pound-mass (\( \text{lbm} \)) at a rate of one foot per second squared (\( \text{ft/s}^2 \)). The equation to define pound-force in terms of pound-mass and acceleration is:
\[
F \, \text{[lbf]} = \frac{m \, \text{[lbm]}}{32.2} \times a \, \text{[ft/s}^2] = m \, \text{[slug]} \times a \, \text{[ft/s}^2]
\]

\subsubsection*{Slug [Mass]}
The slug is yet another English unit, designed to represent mass. It's defined as the amount of mass that will accelerate at \( 1 \, \text{ft/s}^2 \) when acted upon by a force of \( 1 \, \text{lbf} \). The mathematical representation for a slug is:
\[
m \, \text{[slug]} = \frac{F \, \text{[lbf]}}{a \, \text{[ft/s}^2]}
\]

\subsubsection*{Pound-Mass (lbm)}
The pound-mass (\( \text{lbm} \)) is an English unit designated to represent mass in an equivalency with the [lbf]. Although similar to the slug, the pound-mass (\( \text{lbm} \)) has a conversion factor to allow the lbf and lbm to be equivalent in earth's gravity. The mathematical representation is:
\[
m \, \text{[lbm]} = 32.2 \times \frac{F \, \text{[lbf]}}{ a \, \text{[ft/s}^2]}
\]




\newpage
\section*{Examples}
\subsubsection*{Example 1: Convert 1 lbf to lbm and Slug in Earth's Gravity}
One of the most frequent mistakes made by students is not paying close attention to units, particularly in conversions. In the English system, \(1 \, \text{lbm}\) is explicitly defined to have a weight of \(1 \, \text{lbf}\) under Earth's gravity, which accelerates objects at \(32.2 \, \text{ft/s}^2\).
\newline\newline
To demonstrate the equivalence of \(1 \, \text{lbm}\) and \(1 \, \text{lbf}\):
\[\frac{F}{g} = \frac{1 \, \text{[lbf]}}{32.2 \, \text{[ft/s}^2]} \, \times \text{32.2 } = 1 \, \text{[lbm]} \]
Here, it's important to note that by design \(1 \, \text{lbm}\) exerts a force equal to \(1 \, \text{lbf}\) under the influence of Earth's gravity (\(32.2 \, \text{[ft/s}^2]\)), thus why the conversion factor of 32.2 is present. In essence, the units \( \text{[lbm]} \) and \( \text{[lbf]} \) are uniquely equated under Earth's gravitational conditions.
\newline\newline
To convert \(1 \, \text{lbm}\) to \( \text{slug}\):
\[\frac{F}{g} = \frac{1 \, \text{[lbf]}}{32.2 \, \text{[ft/s}^2]} \approx 0.0311 \, \text{[slug]} \]
\subsubsection*{Example 2a: Pressure at the Bottom of a 2-ft High Column of Liquid $[100\text{ lbm}/\text{ft}^3]$}
For the first part, we're looking for the pressure at the bottom of the column in units of \([ \text{lbf/ft}^2 ]\). Starting from the equation \( P = h \times \rho \times g \):
\[ P = 2 \, \text{[ft]} \times 100 \, \text{[lbm/ft}^3] \times 32.2 \, \text{[ft/s}^2] \times \frac{1}{32.2}  = 200 \, \text{[lbf/ft}^2]\]
\subsubsection*{Example 2b: Pressure at the Bottom of a 2-ft High Column of Liquid $[100\text{ slug}/\text{ft}^3]$}
Next, we want to find the pressure given \( \text{[slug/ft}^3] \).
\[ P = 2 \, \text{[ft]} \times 100 \, \text{[slug/ft}^3] \times 32.2 \, \text{[ft/s}^2] = 6440 \, \text{[lbf/ft}^2] \]

\newpage
\subsubsection*{Example 3: Force of an Object Given 100 lbm and 100 Slugs in Different Gravities}

One of the common mistakes is forgetting to adjust for different gravitational constants when converting between mass and force.
\newline \newline
\textbf{In Earth's Gravity (\( 32.2 \, \text{[ft/s}^2] \))}
\newline
1. Calculating mass in \([lbm]\) given \( 100 \, \text{[lbf]} \):
\[
\text{Mass in lbm} = \frac{100 \, \text{[lbf]}}{32.2 \, \text{[ft/s}^2]} \times 32.2 = 100 \, \text{[lbm]} \
\]
Here, \([ft/s^2]\) cancels out in the \([lbf]\), leaving us with \( [lbm] \). 
\newline
\newline
2. Calculating mass in \([slug]\) given \( 100 \, \text{[lbf]} \):
\[
\text{Mass in lbm} = \, \frac{100 \, \text{[lbf]}}{32.2 \, \text{[ft/s}^2]} \approx 3.105 \, \text{[slug]} 
\]
Again, the units \([lbf] / [ft/s^2] = [slug]\).
\newline\newline
\textbf{In a Different Gravity (\( 28 \, \text{[ft/s}^2] \))}
\newline
3. Calculating mass in \([lbm]\) given \( 100 \, \text{[lbf]} \):
\[
\text{Mass in lbm} = \frac{100 \, \text{[lbf]}}{28 \, \text{[ft/s}^2]} \times 32.2 \approx 114.93 \, \text{[lbm]} 
\]
Notice that \([ft/s^2]\) cancels out, leaving us with \( [lbm] \).
\newline\newline
4. Calculating mass in \([slug]\) given \( 100 \, \text{[lbf]} \):
\[
\text{Mass in lbm} = \frac{100 \, \text{[lbf]}}{28 \, \text{[ft/s}^2]} \approx 3.57 \, \text{[slug]}
\]
Once more, the units \([lbf] / [ft/s^2] = [slug]\).


\section*{Common Errors}
A common misconception is using \(\text{lbm}\) when the actual units should be \(\text{slug}\). People often forget to include the factor of \(32.2\) when converting \(\text{lbm}\) to \(\text{slug}\) or vice versa. This mistake can lead to significant errors in calculations, especially those involving forces and pressures.

\end{document}
